\scanpage{130}
\uline{Corollaire}\quad Si $A\in\mathcal L^1(\mathcal H)$, soit
$M_A(r)=\sup_{|z|=r}\det(1+zA)$
\deleted{si $A$ est hermitien positif, et si $A\in\mathcal L^p(\mathcal H)$}, \deleted{\unsure{tel|on}}
\deleted{\unsure{avec}} \inserted{et} $0<p<1$\,; on a
\[
S_p(A)=p\frac{\sin\pi p}{\pi}\int_0^\infty\frac{\log M_{|A|}(r)}{r^{1+p}}\,dr\tag*{(9)}
\]

Il suffit en effet de poser $f(z)=\det(1+z|A|)$,
les $\alpha_i$ étant les valeurs propres de $|A|$\,; on aura
ici $M(r)=f(r)$.

Pour démontrer (4), on \deleted{\unsure{pe|peut}} peut écrire
\[
S_p(A+B)=S_p(|A+B|)=p\frac{\sin\pi p}{\pi}\int_0^\infty
\frac{\log M_{|A+B|}(r)}{r^{1+p}}\,dr
\]
Or $M_{|A+B|}(r)=\det(1+\unsure{r|v}|A+B|)\leq\det(1+r|A|)\det(1+\unsure{r|v}|B|)$
(formule \unsure{\ldots}) \deleted{\unsure{est donc|soit $M_A$}} soit
\[
M_{|A+B|}(r)\leq M_{|A|}(r)\unsure{\cdot|+}M_{|B|}(r)\tag*{(10)}
\]
d'où
\[
S_p(A+B)\leq\frac{p\sin\pi p}{\pi}\int_0^\infty\frac{\log M_{|A|}(r)}{r^{1+\unsure{p|r}}}\,dr
+\frac{p\sin\pi p}{\pi}\int_0^\infty\frac{\log M_{|B|}(r)}{r^{1+\unsure{p|r}}}\,dr
\]
Appliquant de nouveau (9), on obtient bien la
formule voulue (4).

\deleted{Pour prouver}

2) Pour prouver la formule (5), je ne connais
qu'une méthode analytique. De façon générale,
\unsure{.} si $A$ est un opérateur cp\textsuperscript{t} continu, \unsure{soient|on prend|soit} la
suite des val. propres $(\alpha_i)$, \deleted{\unsure{prises|posons}} et si
$0<p<+\infty$, posons $s_p(A)=\sum|\alpha_i|^p$.
\deleted{On a $\unsure{s_p(A)|S_p(A)}=$ \unsure{Si|$S_p$} évidemment pour tout entier $n>0$ $s_p(A)=s_{p/n}(A^n)$. Supposons}
% archive sheet number: 131
